\chapter{Reading the values}

A few conventions hold for every view; once you know them, you can read any dashboard.

\section{What you see is the exact state}
We show no averages and no smoothing. An averaged value is a value the chip never had, and it hides exactly what is
interesting: an arpeggio changing note every frame, an envelope restarting. If a value flickers, the music makes it
flicker.

\section{Notes and frequencies}
Notes are written in English: C, D, E, F, G, A, B, with \# for sharps. Octave 4 holds the 440~Hz A (A4), middle C
is C4. When the frequency falls between two notes we show the nearest; the frequency in hertz gives you the gap. A
slight detune between two voices playing the same note is often deliberate: it thickens the sound.

\section{Levels}
Levels are relative to the chip's maximum: 0~dB is full volume, $-6$~dB half the amplitude, $-\infty$ silence. The
bars follow the same scale.

\section{Registers}
Registers are hexadecimal, prefixed with \$, as in period documentation: \$D418 on the SID, \$804 on the PC Engine.
We show them as the driver wrote them, not as we interpret them; if you look them up in a datasheet you will find the
same values.
